\begin{abstract}
Exa stores document embeddings as short binary codes, groups them into clusters, and scores
the documents of a few clusters with a lookup table. This report asks whether two extra
index structures can reduce that scoring work: \emph{Bitplanes}, which split a cluster by
one sign bit at a time, and \emph{Backward Walk}, which looks up documents by a short key.
We derive the time, memory and recall of all three methods, use the formulas to choose each
method's best settings, and then measure every setting on one fixed dataset: one million
MS~MARCO passages and 200 queries. Only the index parameters change between runs. The
formulas predict the recorded operation counts exactly and the measured search times within
a few percent. On this dataset the original method is fastest at every recall target, and
the calculations show why: on real embeddings, one sign bit removes only about 15\% of the
documents, while a cluster boundary removes about 98\%. The same formulas also state the
conditions under which Bitplanes or Backward Walk would win.
\end{abstract}

\section{Introduction}
\subsection{Exa's method, end to end}
We build on the search pipeline that Exa describes in its engineering blog~\cite{exa}. That
post explains each technique in detail, so we only name the steps here and refer the reader
to it. A document is embedded, the embedding is shortened by keeping its first values, which Matryoshka training makes possible~\cite{mrl}, and each
of the remaining $d$ values is replaced by one bit: 1 if the value is positive, 0 if it is
negative. This is Exa's \emph{binary quantization}. The binary codes are grouped into $C$
clusters. When a query arrives, the system keeps the query in floating point, selects the $P$
clusters whose centres are closest, and scores every document in those clusters with a lookup
table that handles four bits at a time. The best documents can then be reranked with full
vectors. Figure~\ref{fig:pipeline} shows the sequence; the blue box is the step that this
report studies.

\begin{figure}[H]\centering
\begin{tikzpicture}[node distance=5mm and 5mm]
\node[box,text width=24mm] (text) {Document text};
\node[box,text width=30mm,right=of text] (embed) {Embed, keep the first $d$ values};
\node[box,text width=26mm,right=of embed] (bits) {Keep only the signs: $d$ bits};
\node[box,text width=28mm,right=of bits] (store) {Group into $C$ clusters};
\draw[arr] (text)--(embed);\draw[arr] (embed)--(bits);\draw[arr] (bits)--(store);
\node[box,text width=24mm,below=9mm of text] (q) {Query vector\\(floating point)};
\node[box,text width=30mm,right=of q] (route) {Pick the $P$ closest clusters};
\node[scan,text width=26mm,right=of route] (local) {Score the documents in those clusters};
\node[box,text width=28mm,right=of local] (out) {Return the top $K$;\\rerank if wanted};
\draw[arr] (q)--(route);\draw[arr] (route)--(local);\draw[arr] (local)--(out);
\node[font=\footnotesize\itshape,text=Muted] at ($(text.south)!0.5!(q.north)$) {index time above, query time below};
\end{tikzpicture}
\caption{Exa's pipeline as described in~\cite{exa}. The top row happens once, when documents
are indexed. The bottom row happens for every query. The blue box, scoring the documents of
the opened clusters, is the step that the two new methods try to shorten. We call the
unchanged pipeline \textbf{method A}, the original method.}
\label{fig:pipeline}
\end{figure}

Two numbers describe the scoring step. With $N$ documents in $C$ clusters of roughly equal
size, opening $P$ clusters means scoring about $NP/C$ documents. For one million documents,
4096 clusters and one opened cluster, that is about 244 documents. Opening more clusters finds
more of the true answers but scores more documents. The rest of this report is about that
trade.

\subsection{Idea background}
The query tells us exactly which document code would score best. Take a small query with
five values,
\[
 q=[\,0.7,\ 0.5,\ -0.4,\ -0.6,\ 0.3\,].
\]
A document bit of 1 stands for $+1$ and a bit of 0 stands for $-1$; the score is the sum of
$q_i$ times the document's $\pm1$ at each position. At the first position the query value is
$0.7$, so a document bit of 1 contributes $+0.7$ and a bit of 0 contributes $-0.7$. At the
third position the query value is $-0.4$, so now a bit of 0 is better. Following this rule at
every position gives the pattern \texttt{11001}, and its score is
$0.7+0.5+0.4+0.6+0.3=2.5$. No document can score higher.

Every other code loses some of that score. A document whose bit disagrees with the query
sign at position $i$ loses $2|q_i|$: instead of gaining $|q_i|$ it pays $|q_i|$. If we call
the sum of the query weights at the disagreeing positions the \emph{penalty} $p$, then
\begin{equation}
 S(q,b)=Q-2p,\qquad Q=\sum_i|q_i|,\qquad p=\sum_{i:\ b_i\ne\operatorname{sign}(q_i)}|q_i|.
 \label{eq:score}
\end{equation}
The code \texttt{11000} disagrees only at the last position, so its penalty is $0.3$ and its
score is $2.5-0.6=1.9$. The code \texttt{01001} disagrees only at the first position, which
carries the largest weight; its score is $2.5-1.4=1.1$. Position matters: a mismatch on a
large query value costs more than a mismatch on a small one.

This suggests a question. Method A scores every document in the opened clusters, even the
ones that disagree with the query on its largest weights. \emph{What if the index let us go
straight to the documents whose codes are close to the preferred pattern, and skip the
rest?} Two structures do this in different ways. Bitplanes split the documents by one sign
bit at a time, largest query weight first. Backward Walk stores documents under a short key
and visits the keys in order of increasing penalty. Both keep Exa's scoring rule and its
clusters; they change only the blue box of Figure~\ref{fig:pipeline}.

We use the six documents in Table~\ref{tab:toy} as a running example, together with the
query above and $K=2$ results wanted. Scoring all six documents is what method A does. The
best two are documents 0 and 1, with scores 2.5 and 1.9.

\begin{table}[H]\centering\small
\caption{Six example documents with five sign bits each, scored against
$q=[0.7,0.5,-0.4,-0.6,0.3]$. The penalty adds the query weights at the positions where the
document disagrees with the query sign.}
\label{tab:toy}
\begin{tabular}{@{}clcc@{}}\toprule
ID & Code & Disagreeing positions (weights) & Score $=2.5-2p$\\\midrule
0 & \texttt{11001} & none & 2.5\\
1 & \texttt{11000} & 5 (0.3) & 1.9\\
2 & \texttt{10100} & 2, 3, 5 (0.5, 0.4, 0.3) & 0.1\\
3 & \texttt{01001} & 1 (0.7) & 1.1\\
4 & \texttt{11101} & 3 (0.4) & 1.7\\
5 & \texttt{00110} & all five (2.5) & $-2.5$\\\bottomrule
\end{tabular}
\end{table}

\subsection{Two new methods: Bitplanes and Backward Walk}
\subsubsection{Bitplanes}
Method A stores each document as one row of bits. Bitplanes, which we call \textbf{method B},
additionally store the same bits sideways: one row per bit position, one column per document,
inside each cluster. Figure~\ref{fig:bitplane-index} shows both layouts for the example. The
sideways row for position 1 is \texttt{111010}; it says that documents 0, 1, 2 and 4 have a 1
there. Such a row is called a \emph{bitplane}. With 64-bit words, one bitplane for a cluster
of $n$ documents occupies $\lceil n/64\rceil$ words, and the cluster needs $d$ of them.

\begin{figure}[H]\centering
\begin{tikzpicture}
\node[anchor=west,font=\small\bfseries,text=Ink] at (0,3.2) {Method A stores rows (one document per row)};
\matrix (rows) at (1.9,0.9) [matrix of nodes,nodes={cell},column sep=-\pgflinewidth,row sep=-\pgflinewidth]{
 1&1&0&0&1\\ 1&1&0&0&0\\ 1&0&1&0&0\\ 0&1&0&0&1\\ 1&1&1&0&1\\ 0&0&1&1&0\\};
\foreach \i/\l in {1/0,2/1,3/2,4/3,5/4,6/5}{\node[font=\footnotesize,text=Muted,anchor=east] at (rows-\i-1.west) {ID \l\ };}
\node[font=\footnotesize,text=Muted] at (rows-1-1.north) [above=1pt] {b1};
\node[font=\footnotesize,text=Muted] at (rows-1-2.north) [above=1pt] {b2};
\node[font=\footnotesize,text=Muted] at (rows-1-3.north) [above=1pt] {b3};
\node[font=\footnotesize,text=Muted] at (rows-1-4.north) [above=1pt] {b4};
\node[font=\footnotesize,text=Muted] at (rows-1-5.north) [above=1pt] {b5};
\node[anchor=west,font=\small\bfseries,text=Branch] at (7.2,3.2) {Method B adds bitplanes (one bit position per row)};
\matrix (planes) at (10.4,1.2) [matrix of nodes,nodes={cell},column sep=-\pgflinewidth,row sep=-\pgflinewidth]{
 |[hot]|1&|[hot]|1&|[hot]|1&|[hot]|0&|[hot]|1&|[hot]|0\\ 1&1&0&1&1&0\\ 0&0&1&0&1&1\\ 0&0&0&0&0&1\\ 1&0&0&1&1&0\\};
\foreach \i/\l in {1/b1,2/b2,3/b3,4/b4,5/b5}{\node[font=\footnotesize,text=Muted,anchor=east] at (planes-\i-1.west) {\l\ };}
\foreach \j/\l in {1/0,2/1,3/2,4/3,5/4,6/5}{\node[font=\footnotesize,text=Muted] at (planes-1-\j.north) [above=1pt] {\l};}
\node[font=\footnotesize,text=Muted] at (10.4,2.65) {document ID};
\node[font=\footnotesize,text=Branch] at (10.4,-0.75) {highlighted row: the bitplane for position 1};
\draw[arr,Branch] (5.2,0.9) -- node[above,font=\footnotesize,text=Branch]{same bits,} node[below,font=\footnotesize,text=Branch]{turned sideways} (7.3,0.9);
\end{tikzpicture}
\caption{The two layouts of the same six codes. Method A keeps only the rows on the left.
Method B keeps the rows for scoring and adds the bitplanes on the right for splitting. The
highlighted bitplane belongs to position 1, which has the largest query weight.}
\label{fig:bitplane-index}
\end{figure}

Searching uses the bitplanes to narrow down the candidates before scoring
(Figure~\ref{fig:bitplane-search}). We start with a \emph{mask} that marks every document
in the cluster as a candidate: \texttt{111111}. We take the position with the largest query
weight, position 1 with $0.7$, where the query prefers a 1. One AND between the mask and the
bitplane for position 1 gives the documents that agree, \texttt{111010}, which are documents
0, 1, 2 and 4. The remaining documents, 3 and 5, disagree at a position of weight $0.7$, so
every document in that group carries a penalty of at least $0.7$. We keep that group aside
in a queue, labelled with its penalty, and continue with the agreeing group.

The next largest weight is position 4 ($0.6$, prefers 0). All four remaining documents have
a 0 there, so nothing is removed. Position 2 ($0.5$, prefers 1) separates document 2 from
documents 0, 1 and 4. When a group is small enough, which we call the \emph{leaf size}, we
stop splitting and score its documents with the same lookup table as method A. Scoring
documents 0, 1 and 4 gives $2.5$, $1.9$ and $1.7$; the best two are 2.5 and 1.9.

Now the queue decides whether any stored group needs attention. Equation~\ref{eq:score} gives
each group a ceiling: a group with penalty $p$ cannot contain a score above $Q-2p$. The group
$\{2\}$ has penalty $0.5$ and ceiling $1.5$, and the group $\{3,5\}$ has penalty $0.7$ and
ceiling $1.1$. Both ceilings are below $1.9$, the worse of the two scores we already hold, so
neither group can change the answer. The search ends having scored three documents instead
of six, with exactly the answer that method A returns.

\begin{figure}[H]\centering
\begin{tikzpicture}[node distance=4mm]
\node[box,text width=30mm] (route) at (0,0) {Pick the $P$ closest clusters (as in A)};
\node[branch,text width=44mm] (split) at (5.2,0) {Split the candidate mask by one bitplane at a time, largest query weight first};
\node[box,text width=32mm] (score) at (10.6,0) {Score only the documents left in small groups};
\draw[arr] (route)--(split);\draw[arr] (split)--(score);
\node[branch,text width=34mm] (m0) at (5.2,-1.9) {all six: \texttt{111111}};
\node[branch,text width=34mm] (m1) at (2.2,-3.3) {agree at b1: \texttt{111010}\\IDs 0,1,2,4; penalty 0};
\node[box,text width=30mm] (m1x) at (8.4,-3.3) {disagree: \texttt{000101}\\IDs 3,5; penalty 0.7};
\node[branch,text width=34mm] (m2) at (2.2,-4.7) {agree at b4: same four\\penalty 0};
\node[branch,text width=34mm] (m3) at (2.2,-6.1) {agree at b2: \texttt{110010}\\IDs 0,1,4 $\rightarrow$ score};
\node[box,text width=30mm] (m3x) at (8.4,-6.1) {disagree: \texttt{001000}\\ID 2; penalty 0.5};
\draw[arr] (m0)--(m1);\draw[arr] (m0)--(m1x);\draw[arr] (m1)--(m2);\draw[arr] (m2)--(m3);\draw[arr] (m2)--(m3x);
\node[font=\footnotesize,text=Muted,align=left,anchor=west] at (10.3,-4.7) {ceilings $Q-2p$:\\$\{2\}$: $2.5-1.0=1.5$\\$\{3,5\}$: $2.5-1.4=1.1$\\both below 1.9, so skipped};
\end{tikzpicture}
\caption{Method B on the example. Orange marks the changed step and the groups that the
search follows. Each split keeps the disagreeing group in a queue with its penalty. After
scoring the small group, the queued groups are dropped because their ceilings are below the
second-best score already found.}
\label{fig:bitplane-search}
\end{figure}

\begin{lstlisting}
bitplanes(query, opened_clusters, K, leaf_size, node_budget):
    table = score_table(query); best = top_k(K)
    order = positions sorted by |query value|, largest first
    queue = one full mask per opened cluster, penalty 0
    while queue is not empty:
        group = pop the group with the smallest penalty
        if Q - 2*group.penalty < best.worst_score(): break  # no group can win
        if visited groups == node_budget: break              # optional limit
        if group.count <= leaf_size or no positions remain:
            score every document in group.mask with table; offer to best
        else:
            agree, disagree = split group.mask by the bitplane of the next position
            push agree (same penalty); push disagree (penalty + |query value|)
    return best
\end{lstlisting}

Two settings control method B. The \emph{leaf size} decides when a group is small enough to
score. The \emph{node budget} caps how many groups the search may visit; with no budget, the
ceiling rule alone stops the search, and the answer is then the same as method A's within the
opened clusters. Both children of a split are as wide as the whole cluster: a mask of
$\lceil n/64\rceil$ words stays that wide even when only a few bits remain set. This is the
main cost of method B, and Section~2 counts it.

\subsubsection{Backward Walk}
The second structure, which we call \textbf{method C}, makes the ``go straight to the
preferred pattern'' idea literal. Choose $h$ fixed bit positions and read a document's bits
at those positions as a short \emph{key}. Inside each cluster, sort the documents by key so
that documents with the same key sit next to each other, and keep a small \emph{directory}
that maps each key to the start and length of its row range. Figure~\ref{fig:key-index}
shows this for the example with positions 2 and 3 as the key.

\begin{figure}[H]\centering
\begin{tikzpicture}
\node[anchor=west,font=\small\bfseries,text=Keys] at (0,3.3) {Method C sorts the rows by key and adds a directory};
\matrix (rows) at (2.6,1.0) [matrix of nodes,nodes={cell},column sep=-\pgflinewidth,row sep=-\pgflinewidth]{
 1&|[hot]|0&|[hot]|1&0&0\\ 0&|[hot]|0&|[hot]|1&1&0\\ 1&|[hot]|1&|[hot]|0&0&1\\ 1&|[hot]|1&|[hot]|0&0&0\\ 0&|[hot]|1&|[hot]|0&0&1\\ 1&|[hot]|1&|[hot]|1&0&1\\};
\foreach \i/\l in {1/2,2/5,3/0,4/1,5/3,6/4}{\node[font=\footnotesize,text=Muted,anchor=east] at (rows-\i-1.west) {ID \l\ };}
\foreach \j/\l in {1/b1,2/b2,3/b3,4/b4,5/b5}{\node[font=\footnotesize,text=Muted] at (rows-1-\j.north) [above=1pt] {\l};}
\draw[decorate,decoration={brace,amplitude=3pt},Keys] (rows-2-5.east) ++(2pt,-2.5pt) -- ++(0,5.6mm+5.6mm+5pt) node[midway,right=4pt,font=\footnotesize,text=Keys] {key \texttt{01}};
\draw[decorate,decoration={brace,amplitude=3pt},Keys] (rows-5-5.east) ++(2pt,-2.5pt) -- ++(0,3*5.6mm+5pt) node[midway,right=4pt,font=\footnotesize,text=Keys] {key \texttt{10}};
\draw[decorate,decoration={brace,amplitude=3pt},Keys] (rows-6-5.east) ++(2pt,-2.5pt) -- ++(0,5.6mm+5pt) node[midway,right=4pt,font=\footnotesize,text=Keys] {key \texttt{11}};
\node[align=center,font=\small] at (10.3,1.0) {\begin{tabular}{ccc}\toprule
key & start & count\\\midrule
\texttt{01} & 0 & 2\\
\texttt{10} & 2 & 3\\
\texttt{11} & 5 & 1\\\bottomrule
\end{tabular}};
\node[font=\footnotesize,text=Muted] at (10.3,2.35) {directory (key \texttt{00} is absent: no document has it)};
\end{tikzpicture}
\caption{The index of method C for the example, with bit positions 2 and 3 as the key
(highlighted). Rows are stored in key order, so the documents of one key form one contiguous
range. The directory holds one entry per key that occurs. The row IDs double as the list of
documents, so no separate posting list is needed.}
\label{fig:key-index}
\end{figure}

At query time we form the query's preferred key from the same positions. The query prefers
a 1 at position 2 and a 0 at position 3, so its key is \texttt{10}. A directory lookup returns
documents 0, 1 and 3, and we score them: $2.5$, $1.9$ and $1.1$. If this were enough we would
stop; but the best documents may have a different key, so the search must be able to move on.
It does so by walking away from the preferred key one flipped position at a time, cheapest
flip first, which is why we call it a backward walk (Figure~\ref{fig:key-walk}). Flipping
position 3 (weight $0.4$) gives key \texttt{11} with penalty $0.4$; flipping position 2
(weight $0.5$) gives \texttt{00} with penalty $0.5$; flipping both gives \texttt{01} with
penalty $0.9$. As in method B, a key with penalty $p$ has ceiling $Q-2p$. The next key,
\texttt{11}, has ceiling $2.5-0.8=1.7$, which is below the second-best score $1.9$ that we
already hold. The walk stops after scoring three documents, again with method A's answer.

\begin{figure}[H]\centering
\begin{tikzpicture}[node distance=4mm]
\node[box,text width=30mm] (route) at (0,0) {Pick the $P$ closest clusters (as in A)};
\node[keys,text width=44mm] (walk) at (5.2,0) {Look up keys in order of penalty, starting at the query's own key};
\node[box,text width=32mm] (score) at (10.6,0) {Score only the rows those keys list};
\draw[arr] (route)--(walk);\draw[arr] (walk)--(score);
\node[keys,text width=26mm] (k0) at (0.6,-2.2) {\texttt{10}, penalty 0\\rows 0,1,3 $\rightarrow$ score};
\node[box,text width=26mm] (k1) at (3.9,-2.2) {\texttt{11}, penalty 0.4\\ceiling 1.7};
\node[box,text width=26mm] (k2) at (7.2,-2.2) {\texttt{00}, penalty 0.5\\ceiling 1.5};
\node[box,text width=26mm] (k3) at (10.5,-2.2) {\texttt{01}, penalty 0.9\\ceiling 0.7};
\draw[arr] (k0)--(k1);\draw[arr] (k1)--(k2);\draw[arr] (k2)--(k3);
\node[font=\footnotesize,text=Muted,align=center] at (5.5,-3.4) {best two so far: 2.5 and 1.9. The ceiling 1.7 of the next key is already below 1.9, so the walk stops here.};
\end{tikzpicture}
\caption{Method C on the example. Green marks the changed step. Keys are visited in
increasing penalty; each visit is one directory lookup per opened cluster. The walk stops as
soon as the next key's ceiling cannot beat the $K$th best score.}
\label{fig:key-walk}
\end{figure}

\begin{lstlisting}
backward_walk(query, opened_clusters, K, key_positions):
    table = score_table(query); best = top_k(K)
    keys = every key, ordered by penalty against the query's preferred key
    for key, penalty in keys:
        if Q - 2*penalty < best.worst_score(): break  # no key can win
        for cluster in opened_clusters:
            rows = cluster.directory.lookup(key)        # may be empty
            score each row with table; offer to best
    return best
\end{lstlisting}

Method C has one main setting, the key width $h$. There are $2^h$ possible keys. A wide key
makes each key's row range small but multiplies the number of keys the walk may have to
visit; a narrow key does the opposite. The key positions are fixed when the index is built,
so they cannot follow each query's largest weights the way method B's split order does. The
positions the enumeration visits are ordered by the query weights at those $h$ positions only.

Both new methods return exactly method A's answer within the opened clusters when they run to
completion, because the ceiling rule never drops a group or key that could still hold a
top-$K$ document. They can be faster only if the ceilings let them skip enough documents to
pay for the extra reads. Section~2 counts that work, Section~3 uses the counts to choose each
method's settings on the dataset, and Section~4 measures the result.
